\section*{Chapitre \ref{ch:nw:time-synchronisation} --- Synchronisation temporelle}

\begin{solution}{exo:nw:time-synchronisation:1}
Elle gagne \qty{15}{\micro\second} par seconde~: $100/15 = \qty{6.7}{\second}$ pour dépasser 100 microsecondes, $50\,000/15 = 3\,333$~s, soit environ 56
minutes, pour dépasser 50 millisecondes.
\end{solution}

\begin{solution}{exo:nw:time-synchronisation:2}
$\hat\vartheta = \tfrac12[(5\,300 - 0) - (8\,900 - 6\,300)] = \qty{1\,350}{\nano\second}$~; aller-retour $8\,900 - 1\,000 =
\qty{7\,900}{\nano\second}$.
\end{solution}

\begin{solution}{exo:nw:time-synchronisation:3}
L'asymétrie vaut $40 \times 4.88 = \qty{195.1}{\nano\second}$~; sa moitié, \qty{97.5}{\nano\second}, subsiste dans chaque estimation.
\end{solution}

\begin{solution}{exo:nw:time-synchronisation:4}
Le bruit des horodatages (\qty{5}{\micro\second} en logiciel, quelques nanosecondes en matériel) est faible face aux files d'attente de
plusieurs microsecondes par commutateur et par sens, que l'estimation hérite sous la forme d'une asymétrie aléatoire. L'échange le moins retardé parmi huit est
celui qui a le moins attendu dans les deux sens, donc celui dont l'asymétrie est la plus faible.
\end{solution}

\begin{solution}{exo:nw:time-synchronisation:5}
Résoudre $\tfrac12 \times 10^{-5} t^2 + t = 100\,000$~: $t = (\sqrt{1 + 2} - 1)/10^{-5} = 73\,205$~s, soit environ 20.3 heures.
\end{solution}

\begin{solution}{exo:nw:time-synchronisation:6}
En dessous de quelques secondes (les courbes se croisent entre 2 et 8~s)~: les corrections de l'asservissement font osciller l'horloge disciplinée de quelques nanosecondes d'un échange à l'autre, alors que la fréquence
de l'oscillateur libre reste stable sur de courtes durées. Ce qui compte pour la règle est le décalage, que l'asservissement maintient à quelques nanosecondes~; sur
des durées plus longues, c'est la dérive de l'horloge libre qui domine, et là l'horloge disciplinée est meilleure de plusieurs ordres de grandeur.
\end{solution}

\begin{solution}{exo:nw:time-synchronisation:7}
Avec \qty{50}{\micro\second} de file d'attente par commutateur et par sens~: environ \qty{259}{\micro\second} avec les seuls horodatages matériels et
\qty{111}{\micro\second} avec le filtre de délai minimal, tous deux au-delà de la règle~; \qty{25}{\nano\second} avec les \omterm{def:nw:time-synchronisation:bc}{horloges transparentes}, qui suppriment
la file d'attente quelle que soit sa taille.
\end{solution}

\begin{solution}{exo:nw:time-synchronisation:8}
Les statistiques PTP mesurent le décalage par rapport au maître le long du chemin, l'asymétrie fixe du chemin restant invisible~; elles ne disent rien de
l'erreur propre du maître par rapport à l'UTC, de l'étape carte-vers-hôte, ni de l'endroit où l'horodatage est réellement pris. Les 50~ns ne forment qu'un maillon de la chaîne, et non la
chaîne entière.
\end{solution}

\begin{solution}{pb:nw:time-synchronisation:1}
\begin{enumerate}
\item $30 + 20 + 40 + 50 + 50 + 100 + 20 + 500 + 1\,500 = \qty{2\,310}{\nano\second}$.
\item $\sqrt{2\,518\,300} \approx \qty{1\,587}{\nano\second}$.
\item Le point d'horodatage logiciel et l'étape carte-vers-hôte~: $2\,000/2\,310$, soit 87~\%.
\item $100\,000 - 2\,310 = \qty{97\,690}{\nano\second}$.
\item Pire cas \qty{310}{\nano\second}, somme quadratique $\sqrt{18\,300} \approx \qty{135}{\nano\second}$.
\item La carte horodate le paquet avec sa propre horloge, que PTP discipline directement~; l'horloge de l'hôte n'intervient plus.
\item L'asymétrie résiduelle du chemin, \qty{100}{\nano\second}.
\item Calibrer l'asymétrie pour l'éliminer et retirer une \omterm{def:nw:time-synchronisation:bc}{horloge frontière}.
\item Résoudre $\tfrac12 \times 10^{-5} t^2 + t = 97\,690$~: environ 71\,866~s, soit 20.0 heures.
\item Résoudre $\tfrac12 \times 10^{-3} t^2 + 50\,t = 97\,690$~: environ 1\,917~s, soit 32 minutes.
\item Aucun des deux~: 60 heures représentent trois fois le maintien de la meilleure des deux.
\item Consigner chaque minute passée en maintien, déclencher une alarme immédiatement et basculer sur une deuxième référence indépendante, si elle existe.
\item \textbf{Résultat clé.} Le pire cas de la chaîne est \qty{2\,310}{\nano\second} avec des horodatages logiciels et \qty{310}{\nano\second} avec
  des horodatages sur carte~; le grand maître conserve la marge restante pendant 20.0 heures avec le bon oscillateur et 32 minutes avec le moins bon.
\item Les horodatages doivent résoudre une microseconde~: un horodatage logiciel lu sur une horloge à la microseconde y satisfait, mais son point
  d'application doit aussi être identifié et cohérent.
\item Le maintien est borné, quelle que soit la qualité de l'oscillateur~; une référence indépendante garde la chaîne traçable et permet à chaque référence d'être
  vérifiée par l'autre.
\item La conception documentée, les points exacts d'horodatage, les enregistrements de surveillance continue et une revue annuelle.
\item Le décalage transmis par le satellite et une asymétrie calibrée, mesurée par rapport à un récepteur indépendant ou à une référence portable.
\item Ajouter \qty{10}{\micro\second}, plus que le reste de la chaîne réuni, et rester tout de même dans la règle.
\item La limite américaine de 50 millisecondes est plus de vingt mille fois supérieure au pire cas de la chaîne.
\item Chaque maillon entre l'UTC et le point où l'horodatage est pris, avec son incertitude.
\end{enumerate}
\end{solution}

\begin{solution}{iq:nw:time-synchronisation:1}
$t_2 = t_1 + d_{\mathrm{ms}} + \vartheta$, $t_4 = t_3 - \vartheta + d_{\mathrm{sm}}$, donc
$\tfrac12[(t_2 - t_1) - (t_4 - t_3)] = \vartheta + \tfrac12(d_{\mathrm{ms}} - d_{\mathrm{sm}})$. Elle suppose des délais égaux dans les deux sens~;
l'asymétrie ne peut pas être mesurée par les échanges.
\iqlookfor{la dérivation et l'asymétrie non mesurable.}
\end{solution}

\begin{solution}{iq:nw:time-synchronisation:2}
Une \omterm{def:nw:time-synchronisation:bc}{horloge frontière} termine PTP~: elle est esclave en amont et maître en aval. Une \omterm{def:nw:time-synchronisation:bc}{horloge transparente} relaie PTP et corrige chaque
message de son temps de séjour. Les deux suppriment la file d'attente du commutateur de l'estimation, qui est sinon la plus grande erreur sur
un réseau de trading chargé.
\iqlookfor{la file d'attente comme problème, et les deux manières de la supprimer.}
\end{solution}

\begin{solution}{iq:nw:time-synchronisation:3}
PTP avec horodatages matériels et commutateurs participants atteint la nanoseconde, voire quelques dizaines de nanosecondes, sur un réseau local~; NTP avec
horodatages logiciels atteint la microseconde au mieux, la milliseconde sur une grande zone. NTP suffit pour les hôtes dont les horodatages doivent satisfaire
une règle de la milliseconde, comme une exigence de piste d'audit à 50 millisecondes.
\iqlookfor{horodatages matériels, commutateurs participants, et adéquation de l'outil à la règle.}
\end{solution}

\begin{solution}{iq:nw:time-synchronisation:4}
Les journaux de synchronisation des deux horloges à ce moment-là (décalages, maintien, alarmes), l'endroit où chaque horodatage a été pris (carte ou logiciel, émission ou
réception), toute asymétrie fixe des chemins vers le grand maître, et si l'écart tient dans la somme des deux budgets~; puis
chercher un vrai bogue dans l'ordonnancement.
\iqlookfor{les budgets d'erreur avant les conclusions.}
\end{solution}

\begin{solution}{iq:nw:time-synchronisation:5}
Présenter la chaîne documentée et le budget, les enregistrements de surveillance de ce jour-là (verrouillage du grand maître, décalage et délai de chaque esclave, aucun
maintien ni dépassement), la calibration des asymétries fixes, les points d'horodatage et la dernière revue annuelle.
\iqlookfor{des enregistrements et de la traçabilité, pas un chiffre isolé.}
\end{solution}

\begin{solution}{iq:nw:time-synchronisation:6}
Le grand maître passe en maintien et son décalage croît avec l'erreur résiduelle et le vieillissement de son oscillateur~; les esclaves le suivent. Basculer sur
une référence indépendante s'il en existe une, déclencher une alarme, calculer à partir des caractéristiques de l'oscillateur lorsque le budget s'épuise, consigner toute la
période, et réparer avant ce moment.
\iqlookfor{l'arithmétique du maintien et une deuxième référence.}
\end{solution}
